\makereporttitle{2025}{ianuarie}

\makereportpreamble

În luna ianuarie am reușit să progresez semnificativ atât pe proiectul de cercetare pe zona de haos cuantic (studiul apariției haosului într-un model provenit din fizica nucleară), cât și pe cel de calcul numeric (construirea în paralel a matricelor de precondiționare pentru metode iterative). Am continuat să-mi fac temele și să lucrez exercițiile de la cursurile de master la care am fost.

Pe partea de haos cuantic, am dezvoltat o interfață grafică interactivă (GUI) pe baza bibliotecilor \href{https://juliadynamics.github.io/DynamicalSystems.jl/latest/}{DynamicalSystems.jl} și \href{https://docs.makie.org/stable/}{Makie.jl}, cu ajutorul căreia să pot explora mai ușor spațiul fazelor și secțiunile Poincaré pentru sistemul cuadrupol bosonic.

\begin{figure}[htbp]
    \centering
    \includegraphics[width=0.8\linewidth]{reports/2025/images/1/interactive-quantum-chaos-january-2025}
\end{figure}

Am încercat să reproduc și graficul distribuției exponentului Lyapunov maximal în funcție de energia totală a sistemului, dar nu am reușit încă să obțin aceleași rezultate ca în lucrările publicate anterior. Voi continua să investighez sursa acestor diferențe.

Pe partea de construcție în paralel a matricelor de precondiționare, după mai multe încercări eșuate am reușit să identific un caz în care abordarea propusă funcționează (folosind metoda Richardson precondiționată). Schimbarea matricei de precondiționare în timpul iterațiilor a făcut algoritmul să conveargă mai rapid.

\vspace{0.5em}

\textbf{Observațiile conducătorului științific}:

\vspace{3em}

\makereportpostamble{27 ianuarie 2025}
